% L04_overview.tex -- Lecture 4: Fintech Security and Regulation (Overview Variant)
% Frames: ~28 | Charts: subset of 12 (9 used) | Architecture: INTRO/CORE/CLOSING
% Framework: PMSP (Problem -- Method -- Solution -- Practice)
% Audience: MSc Finance/Business, no coding assumed
\input{../../_shared/preamble.tex}

\subtitle{Navigating Compliance in the Digital Finance Era}

% Short versions for footline
\title[Fintech: Regulation]{Financial Technology (FinTech) -- Lecture 4}
\author[J.\ Osterrieder]{Joerg Osterrieder}

\begin{document}

% =============================================
%   INTRO ZONE
%   Frames 1-4 -- Open, motivate, orient
% =============================================

\section{Introduction}

% ---------------------------------------------------------
% Frame 1: Title Page
% ---------------------------------------------------------
\begin{frame}
  \titlepage
  \bottomnote{Lecture 4 of 7 $\cdot$ Financial Technology (FinTech) $\cdot$ MSc Programme $\cdot$ Spring 2026}
\end{frame}

% ---------------------------------------------------------
% Frame 2: Opening Cartoon
% ---------------------------------------------------------
\begin{frame}{``We Need to Comply -- But With What?''}
  \begin{center}
    \includegraphics[width=0.90\textwidth]{11_opening_cartoon/cartoon.pdf}
  \end{center}
  \bottomnote{As of 2024, financial institutions face an average of 257 regulatory changes per day globally (Thomson Reuters). The compliance burden has grown 500\% since the 2008 financial crisis.}
\end{frame}

% ---------------------------------------------------------
% Frame 3: Learning Objectives
% ---------------------------------------------------------
\begin{frame}{Learning Objectives}
  \begin{enumerate}
    \item \textbf{Describe} the spectrum of regulatory approaches to fintech ---
          from innovation-friendly sandboxes to precautionary prohibition ---
          and explain the forces shaping each stance.
          \textcolor{mlgray}{\small[Understand]}
    \item \textbf{Explain} the three stages of money laundering and how KYC, CDD,
          and digital identity frameworks address each stage in practice.
          \textcolor{mlgray}{\small[Understand]}
    \item \textbf{Apply} a regulatory mapping framework to identify which US and
          global agencies govern specific fintech activities and products.
          \textcolor{mlgray}{\small[Apply]}
    \item \textbf{Analyse} how RegTech solutions --- NLP, machine learning, and
          automated reporting --- reduce compliance costs while improving
          detection accuracy.
          \textcolor{mlgray}{\small[Analyse]}
    \item \textbf{Evaluate} the trade-offs between regulatory sandboxes,
          SupTech, and embedded compliance as strategies for managing
          innovation without sacrificing consumer protection.
          \textcolor{mlgray}{\small[Evaluate]}
  \end{enumerate}
  \vspace{0.4em}
  \textcolor{mlpurple}{\textbf{Bloom's levels covered:}} Understand $\to$ Apply $\to$ Analyse $\to$ Evaluate
  \bottomnote{These objectives map directly to the quiz and workshop assessments for this lecture.}
\end{frame}

% ---------------------------------------------------------
% Frame 4: Bridge from L03
% ---------------------------------------------------------
\begin{frame}{Building on L03 -- From Payments to Regulation}
  \begin{columns}[T]
    \begin{column}{0.54\textwidth}
      \textcolor{mlpurple}{\textbf{Where we left off (L03):}}
      \begin{itemize}
        \item Payment systems are undergoing a shift from batch to real-time settlement
        \item Cross-border payments remain slow, opaque, and expensive
        \item Every payment design allocates costs, risks, and power
      \end{itemize}
      \vspace{0.5em}
      \textcolor{mlblue}{\textbf{Where we go today (L04):}}
      \begin{itemize}
        \item \emph{Who} regulates fintech --- and with what authority?
        \item How do \textbf{AML/KYC} requirements shape what fintechs can build?
        \item Can \textbf{RegTech} make compliance a competitive advantage rather
              than a cost centre?
      \end{itemize}
    \end{column}
    \begin{column}{0.42\textwidth}
      \includegraphics[width=\textwidth]{01_regulatory_approaches/chart.pdf}
    \end{column}
  \end{columns}
  \bottomnote{L03 showed you the infrastructure through which money moves. L04 asks: who sets the rules for that infrastructure --- and how are those rules enforced?}
\end{frame}

% =============================================
%   CORE ZONE -- PROBLEM
%   Frames 5-9 -- Regulatory landscape, AML, KYC
% =============================================

\section{Problem: Why Does Fintech Regulation Matter?}

% ---------------------------------------------------------
% Frame 5: The Regulatory Spectrum
% ---------------------------------------------------------
\begin{frame}{The Regulatory Spectrum -- Innovation vs.\ Precaution}
  \begin{columns}[T]
    \begin{column}{0.50\textwidth}
      \includegraphics[width=\textwidth]{01_regulatory_approaches/chart.pdf}
    \end{column}
    \begin{column}{0.46\textwidth}
      Regulators face a fundamental dilemma: \textcolor{mlpurple}{regulate too early}
      and stifle innovation; \textcolor{mlred}{regulate too late} and risk consumer harm.
      \vspace{0.4em}
      \begin{itemize}
        \item \textbf{Innovation-friendly:} Sandboxes, light-touch licensing,
              activity-based regulation (UK FCA, Singapore MAS)
        \item \textbf{Wait-and-see:} Monitor market development, regulate only
              after systemic risk emerges (US approach pre-2022)
        \item \textbf{Precautionary:} Comprehensive frameworks ex ante,
              licensing before launch (EU MiCA, China)
        \item \textbf{Prohibitionist:} Outright bans on certain activities ---
              crypto trading, unlicensed lending (select jurisdictions)
      \end{itemize}
      \vspace{0.3em}
      \begin{block}{Key Insight}
        \textcolor{mlpurple}{No regulator occupies a single position.} Most mix
        approaches by product: sandbox for payments, precaution for lending,
        prohibition for unregistered securities.
      \end{block}
    \end{column}
  \end{columns}
  \bottomnote{The World Bank identifies 50+ countries operating fintech regulatory sandboxes as of 2024. Jurisdictional competition for fintech headquarters is reshaping global capital flows.}
\end{frame}

% ---------------------------------------------------------
% Frame 6: Why Regulation Exists -- The Market Failures
% ---------------------------------------------------------
\begin{frame}{Why Regulate Fintech? The Underlying Market Failures}
  Financial regulation exists because financial markets exhibit persistent
  \textcolor{mlpurple}{market failures} that unregulated competition cannot solve:

  \vspace{0.5em}
  \begin{columns}[T]
    \begin{column}{0.47\textwidth}
      \textcolor{mlpurple}{\textbf{Information asymmetry:}}
      \begin{itemize}
        \item Consumers cannot evaluate algorithmic credit scoring or robo-advice
              quality before purchase
        \item Fintech platforms know user data; users do not know how their data
              is used or monetised
      \end{itemize}
      \vspace{0.3em}
      \textcolor{mlblue}{\textbf{Systemic risk:}}
      \begin{itemize}
        \item Interconnected fintech lending platforms can propagate losses
              across the financial system
        \item Stablecoin runs (TerraUSD, 2022) demonstrated contagion to
              traditional markets
      \end{itemize}
    \end{column}
    \begin{column}{0.47\textwidth}
      \textcolor{mlgreen}{\textbf{Negative externalities:}}
      \begin{itemize}
        \item Money laundering, terrorist financing, and fraud impose costs on
              society that individual firms do not bear
        \item Data breaches affect all users, not just the breached entity
      \end{itemize}
      \vspace{0.3em}
      \begin{alertblock}{The Fintech-Specific Challenge}
        Traditional regulation assumes \textcolor{mlred}{entity-based}
        oversight (regulate the bank). Fintech operates through
        \textcolor{mlred}{activity-based} models that may not fit any
        single licence category.
      \end{alertblock}
    \end{column}
  \end{columns}
  \bottomnote{The entity-vs-activity debate is central to fintech regulation. Should a payments app be regulated as a bank, a money transmitter, a technology company, or something new?}
\end{frame}

% ---------------------------------------------------------
% Frame 7: Money Laundering -- The Three Stages
% ---------------------------------------------------------
\begin{frame}{Money Laundering -- Placement, Layering, Integration}
  \begin{columns}[T]
    \begin{column}{0.50\textwidth}
      \includegraphics[width=\textwidth]{04_money_laundering_stages/chart.pdf}
    \end{column}
    \begin{column}{0.46\textwidth}
      The FATF framework identifies three stages of money laundering:
      \vspace{0.4em}
      \begin{enumerate}
        \item \textcolor{mlpurple}{\textbf{Placement:}} Illicit cash enters the
              financial system --- via deposits, currency exchange, or purchase
              of monetary instruments
        \item \textcolor{mlblue}{\textbf{Layering:}} Complex transactions
              obscure the audit trail --- wire transfers, shell companies,
              trade-based laundering, crypto mixing
        \item \textcolor{mlgreen}{\textbf{Integration:}} Cleaned funds re-enter
              the legitimate economy as investment, real estate, or business
              revenue
      \end{enumerate}
      \vspace{0.3em}
      \begin{block}{Scale}
        The UN estimates \textcolor{mlpurple}{2--5\% of global GDP}
        (USD~800B--2T) is laundered annually. Less than 1\% is
        intercepted.
      \end{block}
    \end{column}
  \end{columns}
  \bottomnote{FATF (Financial Action Task Force) sets global AML standards. Its ``grey list'' of non-compliant jurisdictions directly affects capital flows and correspondent banking access.}
\end{frame}

% ---------------------------------------------------------
% Frame 8: AML Compliance -- The Institutional Response
% ---------------------------------------------------------
\begin{frame}{AML Compliance -- From Suspicion to Enforcement}
  \begin{columns}[T]
    \begin{column}{0.50\textwidth}
      \includegraphics[width=\textwidth]{02_aml_compliance_flow/chart.pdf}
    \end{column}
    \begin{column}{0.46\textwidth}
      AML compliance operates as a multi-layered detection system:
      \vspace{0.4em}
      \begin{itemize}
        \item \textbf{Transaction monitoring:} Real-time screening against
              rules and thresholds (e.g., transactions above USD~10,000)
        \item \textbf{Sanctions screening:} Every customer and counterparty
              checked against OFAC, EU, and UN sanctions lists
        \item \textbf{Suspicious Activity Reports (SARs):} Filed with
              Financial Intelligence Units when patterns suggest laundering
        \item \textbf{Ongoing due diligence:} Risk profiles updated
              continuously, not just at onboarding
      \end{itemize}
      \vspace{0.3em}
      \begin{alertblock}{The False Positive Problem}
        Legacy rule-based systems generate \textcolor{mlred}{95--98\% false
        positives}. Compliance teams spend most of their time
        investigating alerts that lead nowhere.
      \end{alertblock}
    \end{column}
  \end{columns}
  \bottomnote{US banks filed over 4~million SARs in 2023 (FinCEN). The UK's National Crime Agency received 900,000+. Most are never investigated by law enforcement.}
\end{frame}

% ---------------------------------------------------------
% Frame 9: KYC and Digital Identity
% ---------------------------------------------------------
\begin{frame}{KYC -- Know Your Customer in the Digital Age}
  \begin{columns}[T]
    \begin{column}{0.50\textwidth}
      \includegraphics[width=\textwidth]{03_kyc_process_flow/chart.pdf}
    \end{column}
    \begin{column}{0.46\textwidth}
      KYC is the gateway process that determines who can access financial services:
      \vspace{0.4em}
      \begin{itemize}
        \item \textbf{Customer Identification (CIP):} Verify name, DOB, address,
              and government ID --- the minimum legal requirement
        \item \textbf{Customer Due Diligence (CDD):} Assess risk profile,
              source of funds, expected transaction patterns
        \item \textbf{Enhanced Due Diligence (EDD):} For PEPs, high-risk
              jurisdictions, and complex ownership structures
      \end{itemize}
      \vspace{0.3em}
      \textcolor{mlpurple}{\textbf{Digital KYC innovations:}}
      \begin{itemize}
        \item Video verification and biometric matching
        \item eIDAS-compliant electronic identity in the EU
        \item India's Aadhaar: 1.3B biometric IDs enabling paperless onboarding
      \end{itemize}
    \end{column}
  \end{columns}
  \bottomnote{Traditional KYC onboarding costs USD~30--150 per customer and takes 2--12 weeks. Digital KYC can reduce this to under USD~5 and minutes --- but regulatory acceptance varies.}
\end{frame}

% =============================================
%   CORE ZONE -- METHOD
%   Frames 10-15 -- US regulation, global regulation, RegTech
% =============================================

\section{Method: How Fintech Is Regulated}

% ---------------------------------------------------------
% Frame 10: The US Regulatory Patchwork
% ---------------------------------------------------------
\begin{frame}{The US Regulatory Patchwork}
  \begin{columns}[T]
    \begin{column}{0.50\textwidth}
      \includegraphics[width=\textwidth]{05_us_regulatory_patchwork/chart.pdf}
    \end{column}
    \begin{column}{0.46\textwidth}
      The US has no single fintech regulator. Authority is fragmented across:
      \vspace{0.4em}
      \begin{itemize}
        \item \textbf{OCC:} National bank charters; proposed fintech charter
              (contested by states)
        \item \textbf{SEC:} Securities, investment contracts, token offerings
              (Howey test)
        \item \textbf{CFTC:} Commodities and derivatives; claims jurisdiction
              over Bitcoin as a commodity
        \item \textbf{CFPB:} Consumer protection in lending, payments, and
              financial products
        \item \textbf{FinCEN:} AML/BSA compliance for money services businesses
        \item \textbf{State regulators:} 50 separate money transmitter licences
              required for national payment operations
      \end{itemize}
      \vspace{0.2em}
      \begin{block}{The Result}
        A fintech offering lending, payments, and crypto must navigate
        \textcolor{mlpurple}{5+ federal agencies and up to 50 state regulators}
        simultaneously.
      \end{block}
    \end{column}
  \end{columns}
  \bottomnote{The SEC vs.\ CFTC jurisdictional battle over crypto assets remains unresolved. The lack of a unified framework is cited as the primary reason fintech firms relocate to Singapore or the UK.}
\end{frame}

% ---------------------------------------------------------
% Frame 11: US Regulation in Practice -- Three Cases
% ---------------------------------------------------------
\begin{frame}{US Fintech Regulation in Practice}
  \begin{columns}[T]
    \begin{column}{0.31\textwidth}
      \begin{block}{\centering Payments}
        \textbf{Regulator:} FinCEN + 50 states.

        \vspace{0.3em}
        \textbf{Requirement:} Money Transmitter Licence in each state.

        \vspace{0.3em}
        \textbf{Cost:} USD~2--5M+ and 12--24 months for national coverage.

        \vspace{0.3em}
        \textit{Many fintechs partner with a licensed bank instead.}
      \end{block}
    \end{column}
    \begin{column}{0.31\textwidth}
      \begin{block}{\centering Lending}
        \textbf{Regulator:} CFPB + state lending laws.

        \vspace{0.3em}
        \textbf{Requirement:} Fair lending, TILA, ECOA compliance.

        \vspace{0.3em}
        \textbf{Tension:} Bank-fintech partnerships allow rate exportation --- lending at the bank's home state rate.

        \vspace{0.3em}
        \textit{``Madden fix'' and true lender rules remain contested.}
      \end{block}
    \end{column}
    \begin{column}{0.31\textwidth}
      \begin{block}{\centering Crypto Assets}
        \textbf{Regulator:} SEC (securities) vs.\ CFTC (commodities).

        \vspace{0.3em}
        \textbf{Requirement:} Depends on classification --- token-by-token analysis.

        \vspace{0.3em}
        \textbf{Status:} Regulation-by-enforcement (SEC lawsuits) rather than clear statutory framework.

        \vspace{0.3em}
        \textit{FIT21 Act proposed but not enacted as of 2025.}
      \end{block}
    \end{column}
  \end{columns}
  \bottomnote{The bank-partnership model (``rent-a-charter'') allows fintechs to operate nationally under a partner bank's licence. The OCC and FDIC have issued guidance tightening oversight of these arrangements.}
\end{frame}

% ---------------------------------------------------------
% Frame 12: Global Regulatory Comparison
% ---------------------------------------------------------
\begin{frame}{Global Regulatory Comparison}
  \begin{center}
    \includegraphics[width=0.88\textwidth]{06_global_regulatory_comparison/chart.pdf}
  \end{center}
  \vspace{-0.3em}
  \begin{columns}[T]
    \begin{column}{0.47\textwidth}
      \textcolor{mlpurple}{\textbf{Comprehensive frameworks:}}
      \begin{itemize}
        \item \textbf{EU -- MiCA:} World's first comprehensive crypto-asset
              regulation. Stablecoin reserve requirements, exchange licensing,
              consumer disclosure rules. Effective June 2024
        \item \textbf{Singapore -- MAS:} Payment Services Act covers digital
              payments, crypto, and cross-border transfers under one licence
      \end{itemize}
    \end{column}
    \begin{column}{0.47\textwidth}
      \textcolor{mlblue}{\textbf{Sandbox-led approaches:}}
      \begin{itemize}
        \item \textbf{UK -- FCA:} Pioneer of the regulatory sandbox (2016).
              Over 800 firms tested. ``Same risk, same regulation'' principle
        \item \textbf{UAE -- ADGM/DIFC:} Competing free zones offering
              English common law, sandbox entry, and regional gateway status
        \item \textbf{India -- RBI/SEBI:} Sandbox for payments; restrictive
              on crypto; UPI as state-backed alternative
      \end{itemize}
    \end{column}
  \end{columns}
  \bottomnote{MiCA creates a ``passporting'' regime: a crypto firm licensed in one EU member state can operate across all 27. This mirrors the banking passport introduced by the Second Banking Directive (1989).}
\end{frame}

% ---------------------------------------------------------
% Frame 13: Regulatory Sandboxes -- How They Work
% ---------------------------------------------------------
\begin{frame}{Regulatory Sandboxes -- Controlled Experimentation}
  \begin{columns}[T]
    \begin{column}{0.50\textwidth}
      \includegraphics[width=\textwidth]{08_sandbox_adoption_timeline/chart.pdf}
    \end{column}
    \begin{column}{0.46\textwidth}
      A regulatory sandbox provides a \textcolor{mlpurple}{controlled environment}
      where firms test innovative products under relaxed requirements:
      \vspace{0.4em}
      \begin{itemize}
        \item \textbf{Application:} Firm proposes product, test parameters,
              consumer safeguards, and exit plan
        \item \textbf{Authorisation:} Regulator grants limited licence with
              constraints on customer numbers, transaction volume, or duration
        \item \textbf{Testing:} 6--24 months of live operation with enhanced
              reporting and regulator access to data
        \item \textbf{Exit:} Full licence, modification, or discontinuation
              based on results
      \end{itemize}
      \vspace{0.3em}
      \begin{block}{Effectiveness}
        FCA sandbox firms have a \textcolor{mlpurple}{80\% survival rate} vs.\
        60\% for non-sandbox startups. But critics argue sandboxes favour
        well-connected firms over genuine innovators.
      \end{block}
    \end{column}
  \end{columns}
  \bottomnote{The FCA launched the world's first fintech sandbox in 2016. By 2024, 50+ jurisdictions had adopted the model, including the GFIN (Global Financial Innovation Network) for cross-border sandbox testing.}
\end{frame}

% ---------------------------------------------------------
% Frame 14: The Compliance Cost Burden
% ---------------------------------------------------------
\begin{frame}{The Compliance Cost Burden}
  \begin{center}
    \includegraphics[width=0.88\textwidth]{09_compliance_cost_comparison/chart.pdf}
  \end{center}
  \vspace{-0.3em}
  \begin{columns}[T]
    \begin{column}{0.47\textwidth}
      \textcolor{mlpurple}{\textbf{Cost by institution type:}}
      \begin{itemize}
        \item \textbf{Global bank:} USD~1--2B annually on compliance; 10--15\%
              of operating expenses
        \item \textbf{Mid-size bank:} USD~50--100M; 5--8\% of operating expenses
        \item \textbf{Fintech startup:} USD~1--5M; can represent 20--30\% of
              early-stage revenue
      \end{itemize}
    \end{column}
    \begin{column}{0.47\textwidth}
      \textcolor{mlred}{\textbf{The disproportionate burden:}}
      \begin{itemize}
        \item Compliance cost per customer is \textbf{10--50x higher} for small
              institutions than for global banks
        \item AML compliance alone costs the global banking industry
              USD~214B annually (LexisNexis, 2023)
        \item \textcolor{mlpurple}{Compliance is the single largest barrier to
              entry for fintech startups in regulated markets}
      \end{itemize}
    \end{column}
  \end{columns}
  \bottomnote{Post-crisis regulation (Dodd-Frank, AMLD4/5/6, Basel III) increased compliance headcount at major banks by 40--60\%. Many banks now employ more compliance staff than front-office traders.}
\end{frame}

% ---------------------------------------------------------
% Frame 15: RegTech -- The Technology Solution
% ---------------------------------------------------------
\begin{frame}{RegTech -- Technology Meets Compliance}
  \begin{columns}[T]
    \begin{column}{0.50\textwidth}
      \includegraphics[width=\textwidth]{07_regtech_stack_architecture/chart.pdf}
    \end{column}
    \begin{column}{0.46\textwidth}
      RegTech applies technology to reduce compliance costs and improve accuracy:
      \vspace{0.4em}
      \begin{itemize}
        \item \textbf{Identity verification:} Biometrics, document OCR,
              liveness detection for digital KYC
        \item \textbf{Transaction monitoring:} ML models replacing rule-based
              systems; reduce false positives by 50--70\%
        \item \textbf{Regulatory reporting:} Automated extraction and
              submission of required data to regulators
        \item \textbf{Regulatory change management:} NLP parsing of new
              regulations to identify affected processes
        \item \textbf{Risk assessment:} Real-time scoring of customer and
              transaction risk using network analysis
      \end{itemize}
      \vspace{0.3em}
      \begin{block}{Market Size}
        Global RegTech market: \textcolor{mlpurple}{USD~12B (2023)},
        projected to reach USD~35--45B by 2028 (CAGR~25\%).
      \end{block}
    \end{column}
  \end{columns}
  \bottomnote{RegTech is to compliance what fintech is to banking: using technology to do what humans do, faster, cheaper, and at scale. The difference is that RegTech's customer is the regulated institution.}
\end{frame}

% =============================================
%   CORE ZONE -- SOLUTION
%   Frames 16-20 -- RegTech applications, NLP, SupTech, embedded compliance
% =============================================

\section{Solution: RegTech, SupTech, and Embedded Compliance}

% ---------------------------------------------------------
% Frame 16: NLP for Regulatory Intelligence
% ---------------------------------------------------------
\begin{frame}{NLP for Regulatory Intelligence}
  \begin{columns}[T]
    \begin{column}{0.55\textwidth}
      Financial regulation is written in \textcolor{mlpurple}{natural language} ---
      ambiguous, context-dependent, and voluminous. NLP transforms it into
      machine-actionable rules:
      \vspace{0.4em}
      \begin{itemize}
        \item \textbf{Regulatory change detection:} Crawl 1,000+ sources;
              flag changes relevant to the firm's licence and jurisdiction
        \item \textbf{Obligation extraction:} Parse regulatory text to
              identify ``must'', ``shall'', ``prohibited'' --- map to
              internal controls
        \item \textbf{Cross-jurisdictional mapping:} Align equivalent
              requirements across MiCA, FCA Handbook, and Dodd-Frank to
              reduce duplicated compliance effort
        \item \textbf{Gap analysis:} Compare current policies against new
              requirements; generate remediation tasks automatically
      \end{itemize}
    \end{column}
    \begin{column}{0.41\textwidth}
      \begin{block}{Example: CUBE}
        Ingests regulatory feeds from 180+ jurisdictions. Uses NLP to
        extract 50,000+ regulatory obligations. Maps each obligation to
        the firm's product set. Alerts compliance teams to gaps.
      \end{block}
      \vspace{0.3em}
      \begin{alertblock}{The Limitation}
        NLP can parse \textit{what} a regulation says. It cannot determine
        \textit{how} a regulator will \textcolor{mlred}{interpret and
        enforce} it. Regulatory ambiguity is a feature, not a bug.
      \end{alertblock}
    \end{column}
  \end{columns}
  \bottomnote{Thomson Reuters estimates 257 regulatory updates per day globally. No compliance team can manually monitor this volume. NLP-based regulatory intelligence is becoming a necessity, not a luxury.}
\end{frame}

% ---------------------------------------------------------
% Frame 17: ML for Transaction Monitoring
% ---------------------------------------------------------
\begin{frame}{Machine Learning for Transaction Monitoring}
  Traditional rule-based AML systems check transactions against fixed thresholds
  (e.g., ``flag any transfer above USD~10,000''). Machine learning changes the game:

  \vspace{0.5em}
  \begin{columns}[T]
    \begin{column}{0.47\textwidth}
      \textcolor{mlpurple}{\textbf{Supervised approaches:}}
      \begin{itemize}
        \item Train on historically confirmed SARs
        \item Gradient-boosted models achieve 3--5x improvement in
              true positive rates over rule-based systems
        \item Challenge: limited labelled data --- confirmed money
              laundering cases are rare
      \end{itemize}
      \vspace{0.3em}
      \textcolor{mlblue}{\textbf{Unsupervised approaches:}}
      \begin{itemize}
        \item Anomaly detection identifies unusual patterns without
              labelled examples
        \item Network analysis reveals hidden relationships between
              entities and accounts
        \item Detects novel typologies that rules cannot anticipate
      \end{itemize}
    \end{column}
    \begin{column}{0.47\textwidth}
      \textcolor{mlgreen}{\textbf{Impact metrics:}}
      \begin{itemize}
        \item False positive reduction: 50--70\%
        \item True positive improvement: 2--4x
        \item Alert investigation time: reduced by 30--50\%
        \item Compliance FTE savings: 20--40\%
      \end{itemize}
      \vspace{0.3em}
      \begin{alertblock}{The Explainability Problem}
        Regulators require that AML decisions be \textcolor{mlred}{explainable}.
        A black-box model that flags a transaction cannot satisfy the
        regulatory requirement to document \textit{why} the alert was
        raised. Explainable AI (XAI) is now a compliance requirement.
      \end{alertblock}
    \end{column}
  \end{columns}
  \bottomnote{HSBC deployed an AI-based transaction monitoring system in 2022 that reduced false positives by 60\% while doubling the detection of genuinely suspicious activity.}
\end{frame}

% ---------------------------------------------------------
% Frame 18: SupTech -- The Regulator's Digital Toolkit
% ---------------------------------------------------------
\begin{frame}{SupTech -- When the Regulator Goes Digital}
  RegTech helps \textit{firms} comply. \textcolor{mlpurple}{\textbf{SupTech}}
  (Supervisory Technology) helps \textit{regulators} supervise:

  \vspace{0.5em}
  \begin{columns}[T]
    \begin{column}{0.47\textwidth}
      \textcolor{mlpurple}{\textbf{SupTech applications:}}
      \begin{itemize}
        \item \textbf{Automated reporting:} Firms submit machine-readable
              data; regulators ingest and validate automatically
        \item \textbf{Real-time dashboards:} Supervisors monitor systemic
              risk indicators across the entire regulated sector
        \item \textbf{NLP for misconduct detection:} Scan internal
              communications for compliance red flags
        \item \textbf{Network analysis:} Map interconnections between
              financial institutions to identify systemic risk concentrations
      \end{itemize}
    \end{column}
    \begin{column}{0.47\textwidth}
      \textcolor{mlblue}{\textbf{Pioneers:}}
      \begin{itemize}
        \item \textbf{MAS (Singapore):} Automated data collection from
              all licensed entities; ML-based risk scoring
        \item \textbf{Bank of England:} Digital Regulatory Reporting pilot ---
              regulations expressed as executable code
        \item \textbf{De Nederlandsche Bank:} Uses ML to detect anomalies
              in pension fund reporting
      \end{itemize}
      \vspace{0.3em}
      \begin{block}{The Vision}
        ``\textcolor{mlpurple}{Regulation as code}'' --- regulatory
        requirements expressed in machine-executable form, enabling
        automatic compliance verification.
      \end{block}
    \end{column}
  \end{columns}
  \bottomnote{A 2023 BIS survey found that 80\% of central banks and supervisory authorities are exploring or deploying SupTech tools. Budget constraints remain the primary barrier to adoption.}
\end{frame}

% ---------------------------------------------------------
% Frame 19: RegTech Investment and Market Growth
% ---------------------------------------------------------
\begin{frame}{RegTech Investment -- A Market Taking Off}
  \begin{center}
    \includegraphics[width=0.88\textwidth]{10_regtech_investment_trends/chart.pdf}
  \end{center}
  \vspace{-0.3em}
  \begin{columns}[T]
    \begin{column}{0.47\textwidth}
      \textcolor{mlpurple}{\textbf{Investment drivers:}}
      \begin{itemize}
        \item Post-2008 regulatory volume explosion
        \item Rising enforcement fines (USD~10B+ annually)
        \item Digital banking growth requiring scalable compliance
        \item Cross-border expansion requiring multi-jurisdictional coverage
      \end{itemize}
    \end{column}
    \begin{column}{0.47\textwidth}
      \textcolor{mlblue}{\textbf{Leading segments:}}
      \begin{itemize}
        \item \textbf{Identity/KYC:} 30\% of RegTech funding
        \item \textbf{Transaction monitoring:} 25\%
        \item \textbf{Regulatory reporting:} 20\%
        \item \textbf{Risk management:} 15\%
        \item \textbf{Regulatory intelligence:} 10\%
      \end{itemize}
      \vspace{0.2em}
      \textcolor{mlgray}{\small Notable firms: Chainalysis (crypto compliance),
      ComplyAdvantage (AML), Onfido (identity), Ascent (regulatory change).}
    \end{column}
  \end{columns}
  \bottomnote{Global RegTech funding exceeded USD~18B cumulatively through 2024. The sector has produced 10+ unicorns, led by Chainalysis (USD~8.6B valuation at peak).}
\end{frame}

% ---------------------------------------------------------
% Frame 20: Embedded Compliance -- The End State
% ---------------------------------------------------------
\begin{frame}{Embedded Compliance -- Regulation Built Into the Product}
  \begin{center}
    \Large
    \textcolor{mlteal}{\textbf{``The best compliance is the compliance\\
    the user never sees.''}}
  \end{center}
  \vspace{0.5em}
  \begin{columns}[T]
    \begin{column}{0.47\textwidth}
      \textcolor{mlpurple}{\textbf{What embedded compliance means:}}
      \begin{itemize}
        \item KYC happens during onboarding --- a selfie and a photo of
              an ID card, verified in real time
        \item Transaction limits enforced by the software, not by manual
              review
        \item Regulatory reporting generated automatically from
              transaction data
        \item Sanctions screening runs on every payment without human
              intervention
      \end{itemize}
    \end{column}
    \begin{column}{0.47\textwidth}
      \textcolor{mlblue}{\textbf{The compliance-as-a-service model:}}
      \begin{itemize}
        \item API-first RegTech providers (Alloy, Unit, Synapse) offer
              compliance infrastructure as a service
        \item Fintechs embed compliance via APIs rather than building
              in-house teams
        \item Cost: USD~0.10--1.00 per KYC check vs.\ USD~30--150 manual
      \end{itemize}
      \vspace{0.3em}
      \begin{alertblock}{The Risk}
        Outsourced compliance creates \textcolor{mlred}{concentration risk}.
        If a major compliance-as-a-service provider fails, dozens of
        fintechs simultaneously lose their regulatory infrastructure.
      \end{alertblock}
    \end{column}
  \end{columns}
  \bottomnote{Synapse's collapse in 2024 stranded customer funds across multiple fintech partners, illustrating the systemic risk of concentrated compliance infrastructure.}
\end{frame}

% =============================================
%   CORE ZONE -- PRACTICE
%   Frames 21-24 -- Evaluation, synthesis, tensions
% =============================================

\section{Practice: Evidence and Evaluation}

% ---------------------------------------------------------
% Frame 21: A Regulatory Evaluation Framework
% ---------------------------------------------------------
\begin{frame}{A Regulatory Evaluation Framework}
  \begin{columns}[T]
    \begin{column}{0.55\textwidth}
      Five questions to evaluate any fintech regulatory approach:
      \vspace{0.4em}
      \begin{enumerate}
        \item \textbf{Does it protect consumers effectively?}\\
              Are disclosure requirements adequate? Is there recourse
              for harmed consumers?
        \item \textbf{Does it enable or inhibit innovation?}\\
              Can startups enter the market at reasonable cost and speed?
        \item \textbf{Is it proportionate?}\\
              Are compliance costs scaled to the risk and size of the
              regulated entity?
        \item \textbf{Is it technology-neutral?}\\
              Does it regulate the \textit{activity} rather than the
              \textit{technology}, avoiding premature lock-in?
        \item \textbf{Is it enforceable?}\\
              Can the regulator actually detect and punish non-compliance?
      \end{enumerate}
    \end{column}
    \begin{column}{0.41\textwidth}
      \textbf{Applying the framework:}
      \vspace{0.3em}
      \begin{block}{EU MiCA}
        Consumer protection: strong. Innovation: moderate (high
        compliance cost). Proportionality: partial (same rules for
        small and large issuers). Tech-neutral: yes. Enforceability:
        strong (passporting + national supervisors).
      \end{block}
      \vspace{0.25em}
      \begin{block}{US Patchwork}
        Consumer protection: variable. Innovation: inhibited
        (50-state licensing). Proportionality: poor. Tech-neutral: no
        (entity-based). Enforceability: fragmented.
      \end{block}
    \end{column}
  \end{columns}
  \bottomnote{This framework integrates L01's strategy lens with regulatory analysis. Apply it in Workshop D to evaluate a jurisdiction's fintech regulatory approach.}
\end{frame}

% ---------------------------------------------------------
% Frame 22: The Central Trade-Off -- Innovation vs. Protection
% ---------------------------------------------------------
\begin{frame}{The Central Tension -- Innovation, Protection, Inclusion}
  \begin{center}
    \Large
    \textcolor{mlteal}{\textbf{``Regulation that eliminates all risk\\
    also eliminates all innovation.''}}
  \end{center}
  \vspace{0.5em}
  \begin{columns}[T]
    \begin{column}{0.47\textwidth}
      \textcolor{mlgreen}{\textbf{Innovation-enabling regulation:}}
      \begin{itemize}
        \item Sandboxes allow controlled experimentation
        \item Activity-based licensing avoids entity-type lock-in
        \item Proportional requirements reduce startup barriers
        \item Open banking mandates (PSD2) force incumbents to compete
      \end{itemize}
    \end{column}
    \begin{column}{0.47\textwidth}
      \textcolor{mlred}{\textbf{The cost of getting it wrong:}}
      \begin{itemize}
        \item TerraUSD collapse (2022): USD~40B in losses, zero consumer
              recourse
        \item Wirecard fraud (2020): EUR~1.9B missing; BaFin failed to
              detect despite warnings
        \item P2P lending failures (China, 2018): 5,000+ platforms collapsed;
              millions of retail investors lost savings
      \end{itemize}
    \end{column}
  \end{columns}
  \vspace{0.5em}
  \begin{alertblock}{No Free Lunch}
    There is no regulatory design that simultaneously maximises innovation
    speed, consumer protection, and financial inclusion.
    \textcolor{mlred}{\textbf{A jurisdiction must choose its priorities and
    accept the trade-offs.}}
  \end{alertblock}
  \bottomnote{The inclusion--protection trade-off from L02 reappears here in regulatory form. Regulation is how societies formalise where they stand on that trade-off.}
\end{frame}

% ---------------------------------------------------------
% Frame 23: Regulatory Health Indicators
% ---------------------------------------------------------
\begin{frame}{Evaluating Regulatory Health -- Five Indicators}
  \begin{enumerate}
    \item \textcolor{mlpurple}{\textbf{Regulatory clarity:}} Can a fintech founder
          determine which licences are needed, what they cost, and how long the
          process takes --- \emph{before} engaging lawyers?
    \vspace{0.3em}
    \item \textcolor{mlblue}{\textbf{Proportionality:}} Are compliance requirements
          scaled to firm size and risk? Does a payment startup face the same burden
          as a systemically important bank?
    \vspace{0.3em}
    \item \textcolor{mlgreen}{\textbf{Enforcement consistency:}} Are rules applied
          equally to incumbents and challengers, or do legacy institutions receive
          de facto leniency?
    \vspace{0.3em}
    \item \textcolor{mlorange}{\textbf{International coherence:}} Can a firm licensed
          in one jurisdiction operate across borders without duplicating the entire
          compliance process?
    \vspace{0.3em}
    \item \textcolor{mlteal}{\textbf{Adaptability:}} Can the regulatory framework
          accommodate technologies that did not exist when the rules were written?
  \end{enumerate}
  \vspace{0.3em}
  \begin{block}{}
    These five indicators apply to any jurisdiction's fintech regulatory approach.
    Use them in Workshop D.
  \end{block}
  \bottomnote{Apply these indicators to a jurisdiction you know. Which indicator reveals the sharpest weakness? Bring your analysis to the workshop.}
\end{frame}

% ---------------------------------------------------------
% Frame 24: What Comes Next
% ---------------------------------------------------------
\begin{frame}{What Comes Next}
  \begin{itemize}
    \item \textcolor{mlpurple}{\textbf{Next: L05 (Wealth Management and Robo-Advisory)}} \\
          How algorithmic advice is transforming investment management.
          Portfolio construction, risk profiling, fiduciary duty, and the
          democratisation of wealth management.
    \vspace{0.4em}
    \item \textcolor{mlblue}{\textbf{Before L05, reflect:}}
      \begin{itemize}
        \item When you opened your last bank or fintech account, what KYC steps
              did you go through? How long did it take? Could it have been simpler?
        \item Pick a recent fintech enforcement action (SEC, FCA, or MAS).
              What went wrong and which regulatory gap was exposed?
      \end{itemize}
    \vspace{0.4em}
    \item \textcolor{mlgreen}{\textbf{Workshop D preparation:}}
      Apply the five regulatory health indicators (Frame~23) to one jurisdiction's
      approach to fintech regulation. Bring a two-paragraph evaluation to class.
  \end{itemize}
  \vspace{0.4em}
  \begin{block}{Course Arc}
    L01: Foundations $\to$ L02: Ecosystem $\to$ L03: Payments $\to$
    \textbf{L04: Regulation} $\to$ L05: Wealth $\to$ L06: Insurance $\to$ L07: Technology
  \end{block}
  \bottomnote{All lecture slides and workshop case materials are available on the course website.}
\end{frame}

% =============================================
%   CLOSING ZONE
%   Frames 25-28 -- Cartoon, Takeaways, Summary
% =============================================

\section{Closing}

% ---------------------------------------------------------
% Frame 25: Closing Cartoon
% ---------------------------------------------------------
\begin{frame}{``But Our Algorithm Said It Was Compliant''}
  \begin{center}
    \includegraphics[width=0.90\textwidth]{12_closing_cartoon/cartoon.pdf}
  \end{center}
  \bottomnote{Global AML fines exceeded USD~5~billion in 2023. The largest single penalty: Binance paid USD~4.3B to settle US DOJ and FinCEN charges for systematic AML failures.}
\end{frame}

% ---------------------------------------------------------
% Frame 26: Key Takeaways
% ---------------------------------------------------------
\begin{frame}{Key Takeaways}
  \begin{enumerate}
    \item \textcolor{mlpurple}{\textbf{Regulation exists for a reason:}} Information
          asymmetry, systemic risk, and negative externalities are real market
          failures that unregulated fintech markets cannot self-correct
    \vspace{0.15em}
    \item \textcolor{mlblue}{\textbf{The regulatory spectrum:}} Jurisdictions range
          from innovation-friendly sandboxes to precautionary frameworks. No single
          approach dominates; most regulators mix strategies by product type
    \vspace{0.15em}
    \item \textcolor{mlgreen}{\textbf{AML/KYC is the gateway:}} Anti-money laundering
          and know-your-customer requirements shape what fintechs can build, whom they
          can serve, and at what cost
    \vspace{0.15em}
    \item \textcolor{mlorange}{\textbf{The US patchwork is uniquely fragmented:}}
          Five-plus federal agencies and 50 state regulators create the world's most
          complex compliance environment for fintech
    \vspace{0.15em}
    \item \textcolor{mlteal}{\textbf{MiCA sets the global template:}} The EU's
          comprehensive crypto-asset regulation is being studied and adapted worldwide
    \vspace{0.15em}
    \item \textcolor{mlred}{\textbf{RegTech transforms compliance:}} ML-based
          transaction monitoring, NLP regulatory intelligence, and API-based KYC can
          reduce compliance costs by 30--70\% while improving detection accuracy
    \vspace{0.15em}
    \item \textcolor{mlpurple}{\textbf{Embedded compliance is the future:}}
          Compliance built into the product via APIs rather than bolted on via manual
          processes --- but concentration risk is the new systemic concern
  \end{enumerate}
  \bottomnote{Review question: Apply the five-question regulatory evaluation framework (Frame~21) to a jurisdiction of your choice. What does the framework reveal?}
\end{frame}

% ---------------------------------------------------------
% Frame 27: Summary and Key Vocabulary
% ---------------------------------------------------------
\begin{frame}{Summary and Key Vocabulary}
  \begin{block}{Lecture Summary}
    Fintech regulation is the arena where innovation meets accountability.
    Regulators worldwide are experimenting with sandboxes, activity-based
    licensing, and SupTech to manage the tension between enabling innovation
    and protecting consumers. AML/KYC requirements remain the single most
    consequential regulatory constraint on fintech product design, determining
    who can be served and at what cost. The US regulatory patchwork creates
    unique barriers, while the EU's MiCA framework establishes the first
    comprehensive template for digital asset regulation. RegTech is
    transforming compliance from a cost centre into a potential competitive
    advantage, but the central lesson remains:
    \textcolor{mlpurple}{\textbf{regulatory design is not a technical
    problem --- it is a societal choice about how much risk a jurisdiction
    is willing to accept in exchange for innovation and inclusion.}}
  \end{block}
  \vspace{0.4em}
  \begin{multicols}{2}
    \small
    \begin{itemize}
      \item \textbf{Regulatory Sandbox}
      \item \textbf{AML / KYC / CDD / EDD}
      \item \textbf{FATF Grey List}
      \item \textbf{Money Laundering Stages}
      \item \textbf{MiCA (Markets in Crypto-Assets)}
      \item \textbf{OCC / SEC / CFTC / CFPB}
      \item \textbf{RegTech / SupTech}
      \item \textbf{Suspicious Activity Report (SAR)}
      \item \textbf{Compliance-as-a-Service}
      \item \textbf{Regulation as Code}
    \end{itemize}
  \end{multicols}
  \vspace{0.2em}
  \textcolor{mlgray}{\small \textbf{Next lecture:} Wealth Management and Robo-Advisory.
  Algorithmic portfolio construction, risk profiling, fiduciary duty, and the
  democratisation of investment management.}
  \bottomnote{Bring your Workshop D regulatory evaluation to the L05 opening discussion.}
\end{frame}

\end{document}
